\documentclass{article}

\begin{document}
	
	SUPPORT VECTOR MACHINES — COMPLETE EXAM WITH ANSWERS
	
	\bigskip
	==================================================
	
	PARAGRAPH 1 — INTRODUCTION TO SVM
	
	\begin{enumerate}
		\item What is the primary task of a Support Vector Machine?
		\begin{enumerate}
			\item Clustering
			\item Dimensionality reduction
			\item Classification
			\item Reinforcement learning
		\end{enumerate}
		
		\item SVMs originate from:
		\begin{enumerate}
			\item Bayesian inference
			\item Statistical Learning Theory
			\item Information theory
			\item Graph theory
		\end{enumerate}
		
		\item Basic SVMs are mainly designed for:
		\begin{enumerate}
			\item Regression
			\item Binary classification
			\item Multilabel classification
			\item Unsupervised learning
		\end{enumerate}
		
		\item A key advantage of SVMs is:
		\begin{enumerate}
			\item No hyperparameters
			\item Strong theoretical foundation
			\item Automatic feature learning
			\item Guaranteed zero error
		\end{enumerate}
		
		\item The output of a standard SVM is:
		\begin{enumerate}
			\item Probability
			\item Continuous score only
			\item Binary label
			\item Cluster index
		\end{enumerate}
	\end{enumerate}
	
	Correct answers: 1C, 2B, 3B, 4B, 5C
	
	\bigskip
	==================================================
	
	PARAGRAPH 2 — TRAINING DATA
	
	\begin{enumerate}
		\item Each SVM training example consists of:
		\begin{enumerate}
			\item Two feature vectors
			\item One feature vector and one label
			\item One label only
			\item Two labels
		\end{enumerate}
		
		\item Feature vectors belong to:
		\begin{enumerate}
			\item Real numbers
			\item Two-dimensional space
			\item n-dimensional space
			\item Natural numbers
		\end{enumerate}
		
		\item Class labels are usually:
		\begin{enumerate}
			\item Continuous
			\item Zero and one
			\item Minus one and plus one
			\item One, two, three
		\end{enumerate}
		
		\item The index i refers to:
		\begin{enumerate}
			\item Feature index
			\item Dimension index
			\item Sample index
			\item Class index
		\end{enumerate}
		
		\item SVM training data must be:
		\begin{enumerate}
			\item Unlabeled
			\item Weakly labeled
			\item Fully labeled
			\item Sequential
		\end{enumerate}
	\end{enumerate}
	
	Correct answers: 1B, 2C, 3C, 4C, 5C
	
	\bigskip
	==================================================
	
	PARAGRAPH 3 — HYPERPLANE AND DECISION FUNCTION
	
	\begin{enumerate}
		\item The separating hyperplane is defined by:
		\begin{enumerate}
			\item w + x + b = 0
			\item w transpose x plus b equals zero
			\item x transpose x plus b equals zero
			\item w transpose w equals one
		\end{enumerate}
		
		\item The vector w controls:
		\begin{enumerate}
			\item Orientation of the hyperplane
			\item Offset of the hyperplane
			\item Margin width
			\item Kernel choice
		\end{enumerate}
		
		\item The bias term b controls:
		\begin{enumerate}
			\item Orientation
			\item Offset
			\item Dimensionality
			\item Regularization
		\end{enumerate}
		
		\item A point is on the decision boundary if:
		\begin{enumerate}
			\item w transpose x plus b is positive
			\item w transpose x plus b is negative
			\item w transpose x plus b equals zero
			\item w transpose x plus b equals one
		\end{enumerate}
		
		\item Classification is based on:
		\begin{enumerate}
			\item Distance from origin
			\item Magnitude of x
			\item Sign of w transpose x plus b
			\item Kernel value only
		\end{enumerate}
	\end{enumerate}
	
	Correct answers: 1B, 2A, 3B, 4C, 5C
	
	\bigskip
	==================================================
	
	PARAGRAPH 4 — MARGIN AND SUPPORT VECTORS
	
	\begin{enumerate}
		\item The margin is:
		\begin{enumerate}
			\item Distance between centroids
			\item Distance between margin hyperplanes
			\item Distance from origin
			\item Distance between all samples
		\end{enumerate}
		
		\item SVM optimization aims to:
		\begin{enumerate}
			\item Minimize training error only
			\item Maximize likelihood
			\item Maximize margin
			\item Maximize variance
		\end{enumerate}
		
		\item Support vectors are:
		\begin{enumerate}
			\item All samples
			\item Misclassified samples
			\item Samples on the margin
			\item Random samples
		\end{enumerate}
		
		\item Which samples define the classifier?
		\begin{enumerate}
			\item All samples
			\item Support vectors
			\item Nearest neighbors
			\item Class centroids
		\end{enumerate}
		
		\item Larger margins usually:
		\begin{enumerate}
			\item Increase overfitting
			\item Reduce robustness
			\item Improve generalization
			\item Have no effect
		\end{enumerate}
	\end{enumerate}
	
	Correct answers: 1B, 2C, 3C, 4B, 5C
	
	\bigskip
	==================================================
	
	PARAGRAPH 5 — SOFT MARGIN AND PARAMETER C
	
	\begin{enumerate}
		\item Slack variables are introduced to:
		\begin{enumerate}
			\item Increase margin
			\item Allow errors
			\item Remove kernels
			\item Normalize data
		\end{enumerate}
		
		\item The parameter C controls:
		\begin{enumerate}
			\item Kernel type
			\item Margin–error tradeoff
			\item Dimensionality
			\item Feature scaling
		\end{enumerate}
		
		\item A large value of C emphasizes:
		\begin{enumerate}
			\item Wider margin
			\item Fewer errors
			\item More regularization
			\item Lower complexity
		\end{enumerate}
		
		\item A small value of C results in:
		\begin{enumerate}
			\item Narrow margin
			\item Hard-margin behavior
			\item Wider margin with more errors
			\item No solution
		\end{enumerate}
		
		\item Soft-margin SVM is mainly used to:
		\begin{enumerate}
			\item Handle noisy data
			\item Eliminate support vectors
			\item Avoid kernels
			\item Reduce dimensions
		\end{enumerate}
	\end{enumerate}
	
	Correct answers: 1B, 2B, 3B, 4C, 5A
	
	\bigskip
	==================================================
	
	PARAGRAPH 6 — NONLINEAR SVM AND KERNELS
	
	\begin{enumerate}
		\item Kernels are used to:
		\begin{enumerate}
			\item Reduce dimensionality
			\item Compute dot products in feature space
			\item Remove slack variables
			\item Normalize data
		\end{enumerate}
		
		\item The kernel trick avoids:
		\begin{enumerate}
			\item Optimization
			\item Explicit feature mapping
			\item Regularization
			\item Classification
		\end{enumerate}
		
		\item Which kernel corresponds to infinite dimensions?
		\begin{enumerate}
			\item Linear
			\item Polynomial
			\item Gaussian (RBF)
			\item Sigmoid
		\end{enumerate}
		
		\item A linear kernel is equivalent to:
		\begin{enumerate}
			\item Nonlinear SVM
			\item Hard-margin SVM always
			\item Linear SVM
			\item Neural network
		\end{enumerate}
		
		\item Kernel choice mainly affects:
		\begin{enumerate}
			\item Data labels
			\item Feature space geometry
			\item Number of samples
			\item Training size
		\end{enumerate}
	\end{enumerate}
	
	Correct answers: 1B, 2B, 3C, 4C, 5B
	
	\bigskip
	==================================================
	
	PARAGRAPH 7 — MULTICLASS SVM
	
	\begin{enumerate}
		\item Multiclass SVMs are usually implemented using:
		\begin{enumerate}
			\item Softmax
			\item One-vs-all
			\item K-means
			\item PCA
		\end{enumerate}
		
		\item In one-vs-all with K classes, how many classifiers are trained?
		\begin{enumerate}
			\item K
			\item K-1
			\item K(K-1)
			\item 2K
		\end{enumerate}
		
		\item Each binary classifier separates:
		\begin{enumerate}
			\item One class vs all others
			\item All classes simultaneously
			\item Pairs of classes
			\item Random subsets
		\end{enumerate}
		
		\item Final class decision is based on:
		\begin{enumerate}
			\item Random choice
			\item Highest decision score
			\item Smallest margin
			\item Kernel value
		\end{enumerate}
		
		\item One-vs-all strategy is:
		\begin{enumerate}
			\item Exact
			\item Heuristic but effective
			\item Unsupervised
			\item Kernel-free
		\end{enumerate}
	\end{enumerate}
	
	Correct answers: 1B, 2A, 3A, 4B, 5B
	
\end{document}
